
\documentclass[12pt,a4paper,oneside,parskip=half]{scrreprt}

%--------------------------------------------------------------------
% Sprache, Kodierung und Seitenlayout
%--------------------------------------------------------------------
\usepackage[utf8]{inputenc}
\usepackage[T1]{fontenc}
\usepackage[ngerman]{babel}
\usepackage[autostyle=true,german=quotes]{csquotes}
\usepackage{lmodern}
\usepackage{microtype}
\usepackage[
    left=3.5cm,
    right=2.5cm,
    top=2.5cm,
    bottom=2.5cm,
    headheight=18pt,
    headsep=6mm,
    footskip=12mm
]{geometry}
\usepackage{setspace}
\usepackage{scrlayer-scrpage}
\usepackage{xurl}
\usepackage{hyperref}
\usepackage{bookmark}
%--------------------------------------------------------------------
% Tabellen, Listen, Abbildungen und Code
%--------------------------------------------------------------------
\usepackage{booktabs}
\usepackage{longtable}
\usepackage{array}
\usepackage{tabularx}
\usepackage{enumitem}
\usepackage{listings}
\usepackage{xcolor}
\usepackage{float}
\usepackage{caption}
\usepackage{graphicx}

\onehalfspacing
\setcounter{tocdepth}{2}
\setcounter{secnumdepth}{3}

\clearpairofpagestyles
\ofoot*{\pagemark}
\pagestyle{scrheadings}
\renewcommand*{\titlepagestyle}{empty}
\renewcommand*{\chapterpagestyle}{scrheadings}
\setkomafont{pagefoot}{\normalfont\small}
%--------------------------------------------------------------------
% Einheitliche Schrift im gesamten Dokument
%--------------------------------------------------------------------
\setkomafont{disposition}{\normalfont\bfseries}
\setkomafont{chapter}{\normalfont\huge\bfseries}
\setkomafont{section}{\normalfont\Large\bfseries}
\setkomafont{subsection}{\normalfont\large\bfseries}
\setkomafont{subsubsection}{\normalfont\normalsize\bfseries}

\setkomafont{pageheadfoot}{\normalfont\small}
\setkomafont{pagefoot}{\normalfont\small}
\setkomafont{caption}{\normalfont\small}
\setkomafont{captionlabel}{\normalfont\small\bfseries}
\captionsetup[table]{justification=centering,singlelinecheck=false}
\renewcommand{\arraystretch}{1.25}
\setlength{\arrayrulewidth}{0.4pt}
\setlength{\tabcolsep}{6pt}
\setlength{\extrarowheight}{2pt}
\setkomafont{descriptionlabel}{\normalfont\bfseries}
\setkomafont{title}{\normalfont\bfseries}
\setkomafont{author}{\normalfont}
\setkomafont{date}{\normalfont}
%--------------------------------------------------------------------
% Einheitliche obere Position von Kapitelanfängen
%--------------------------------------------------------------------
\RedeclareSectionCommand[
    beforeskip=0pt,
    afterskip=0.8\baselineskip
]{chapter}

\hypersetup{
    colorlinks=true,
    linkcolor=black,
    citecolor=black,
    urlcolor=black,
    unicode=true,
    pdfauthor={Rashed Alsuhaibi},
    pdftitle={Hostbasierte Erkennung obfuskierter PowerShell-Aktivität mittels ScriptBlock-Logging und Wazuh}
}

\lstdefinestyle{code}{
    basicstyle=\ttfamily\small,
    breaklines=true,
    frame=single,
    columns=fullflexible,
    keepspaces=true,
    showstringspaces=false,
    numbers=left,
    numberstyle=\tiny,
    xleftmargin=2em,
    framexleftmargin=1.5em,
    literate={ä}{{\"a}}1 {ö}{{\"o}}1 {ü}{{\"u}}1 {Ä}{{\"A}}1 {Ö}{{\"O}}1 {Ü}{{\"U}}1 {ß}{{\ss}}1
}
\lstset{style=code}
\renewcommand{\lstlistingname}{Codeblock}
\renewcommand{\lstlistlistingname}{Quell- und Pseudocodeverzeichnis}

\newcommand{\platzhalter}[1]{\textcolor{gray}{\textbf{[Einfügen: #1]}}}
\newcommand{\metric}[1]{\textit{#1}}

\newcolumntype{Y}{>{\raggedright\arraybackslash}X}
\newcolumntype{L}[1]{>{\raggedright\arraybackslash}p{#1}}
\newcolumntype{C}[1]{>{\centering\arraybackslash}p{#1}}

\newcommand{\codecell}[1]{\begingroup\Urlmuskip=0mu plus 1mu\relax\nolinkurl{#1}\endgroup}
\newcommand{\codebreak}[2]{\codecell{#1#2}}
\newcommand{\codebreakthree}[3]{\codecell{#1#2#3}}

\newcommand{\screenshotplaceholder}[3]{%
\begin{figure}[H]
\centering
\IfFileExists{figures/#1}{%
    \includegraphics[width=0.9\textwidth]{figures/#1}%
}{%
    \fbox{%
    \begin{minipage}{0.88\textwidth}
    \centering
    \vspace{0.8cm}
    \textbf{Screenshot/Abbildung einfügen}\\[0.3cm]
    Dateiname: \path{figures/#1}\\[0.2cm]
    Inhalt: #2
    \vspace{0.8cm}
    \end{minipage}}
}
\caption{#2}
\label{#3}
\end{figure}
}

% Title page info
\title{Hostbasierte Erkennung obfuskierter PowerShell-Aktivität mittels ScriptBlock-Logging und Wazuh \\[1ex] \large
Ein Vergleich regelbasierter Wazuh-Erkennung und lexikalischer Merkmale auf Basis von ScriptBlock-Logging unter Windows 11}
\author{Rashed Alsuhaibi}
\date{\today}

\begin{document}

%--------------------------------------------------------------------
% Titelseite und Verzeichnisse
%--------------------------------------------------------------------
\pagenumbering{roman}
\maketitle

\chapter*{Eidesstattliche Erklärung}
\addcontentsline{toc}{chapter}{Eidesstattliche Erklärung}
Hiermit versichere ich, dass ich die vorliegende Arbeit selbstständig und ohne Benutzung anderer als der angegebenen Hilfsmittel angefertigt habe. Alle Stellen, die wörtlich oder sinngemäß aus veröffentlichten und nicht veröffentlichten Quellen entnommen wurden, sind als solche kenntlich gemacht. Die Arbeit wurde in gleicher oder ähnlicher Form noch keiner anderen Prüfungsbehörde vorgelegt.

\vspace{2cm}
\begin{tabular}{p{6cm}p{6cm}}
Ort, Datum & Unterschrift\\[1.2cm]
\rule{6cm}{0.4pt} & \rule{6cm}{0.4pt}
\end{tabular}
\clearpage

\chapter*{Danksagung}
\addcontentsline{toc}{chapter}{Danksagung}
Ich möchte mich bei allen Personen bedanken, die mich während der Erstellung dieser Bachelorarbeit fachlich und persönlich unterstützt haben. Mein besonderer Dank gilt der betreuenden Person für die fachliche Begleitung und die Rückmeldungen zur Eingrenzung des Themas. Ebenfalls danke ich Familie, Freunden und Kommilitoninnen und Kommilitonen für Motivation, Unterstützung und konstruktive Gespräche während der Bearbeitungszeit.
\clearpage

\chapter*{Zusammenfassung}
\addcontentsline{toc}{chapter}{Zusammenfassung}
Diese Bachelorarbeit untersucht die hostbasierte Erkennung obfuskierter PowerShell-Aktivität unter Windows 11. Im Mittelpunkt steht die Frage, wie gut standardmäßige Wazuh-Regeln verdächtige PowerShell-ScriptBlocks erkennen und ob eine einfache lexikalische Analyse diese regelbasierte Erkennung sinnvoll ergänzen kann. Dazu wird eine kontrollierte Laborumgebung mit aktiviertem ScriptBlock-Logging, einem Wazuh-Agenten und einer zentralen Wazuh-Instanz aufgebaut.

Für die Evaluation wurde ein kuratierter Datensatz in der Version 4 erstellt. Dieser umfasst 300 PowerShell-ScriptBlock-Samples, davon 150 legitime und 150 simuliert maliziöse Samples. Die legitimen Samples bestehen aus 100 normalen administrativen und 50 bewusst schwerer abzugrenzenden benignen Fällen. Die simuliert maliziösen Samples bestehen aus 30 direkten Indikator-Samples und 120 obfuskierten Varianten. Betrachtet werden Base64-Indikatoren, String-Konkatenation, Backticks, gemischte Groß- und Kleinschreibung, Alias- beziehungsweise Befehlsrekonstruktion sowie kombinierte Obfuskation. Alle simuliert maliziösen Samples sind ungefährliche, nicht-ausführende Indikator-Samples.

Die Eigenleistung besteht aus dem Aufbau des Evaluierungsdesigns, der Kuratierung des Datensatzes, der Konzeption der Wazuh-Standardregel-Baseline sowie der Implementierung einer Python-Pipeline zur Feature-Extraktion, lexikalischer Bewertung, Metrikberechnung und Zuordnung exportierter Wazuh-Alerts zu einzelnen Datensatz-Samples. Im finalen Laborlauf wurden 69 relevante Wazuh-Alert-Zeilen exportiert und 64 eindeutigen Samples zugeordnet. Die Wazuh-Standardregeln erreichten eine Precision von 0{,}4688, einen Recall von 0{,}2000, einen F1-Score von 0{,}2804 und eine FPR von 0{,}2267. Die lexikalische Pipeline erreichte bei Schwellenwert 6 eine Precision von 0{,}9015, einen Recall von 0{,}7933, einen F1-Score von 0{,}8440 und eine FPR von 0{,}0867.
\clearpage

\chapter*{Abstract}
\addcontentsline{toc}{chapter}{Abstract}
This bachelor thesis investigates host-based detection of obfuscated PowerShell activity on Windows 11. The main objective is to evaluate how effectively standard Wazuh rules detect suspicious PowerShell script blocks and whether a simple lexical analysis can complement rule-based detection. For this purpose, a controlled laboratory environment is created using PowerShell Script Block Logging, a Wazuh agent, and a central Wazuh instance.

A curated version 4 dataset is prepared for the evaluation. It contains 300 PowerShell ScriptBlock samples, consisting of 150 benign and 150 simulated malicious samples. The benign class contains 100 normal administrative samples and 50 hard-benign cases. The simulated malicious class contains 30 direct indicator samples and 120 obfuscated variants. The dataset covers Base64 indicators, string concatenation, backtick literals, mixed casing, alias or command reconstruction, and combined obfuscation. All simulated malicious samples are safe, non-executable indicator samples.

The practical contribution includes the evaluation design, dataset curation, the Wazuh standard-rule baseline concept, and the implementation of a Python pipeline for feature extraction, lexical scoring, metric calculation, and mapping exported Wazuh alerts back to individual dataset samples. In the final laboratory run, 69 relevant Wazuh alert lines were exported and mapped to 64 unique samples. The Wazuh standard rules achieved a precision of 0.4688, recall of 0.2000, F1-score of 0.2804, and false-positive rate of 0.2267. Using threshold 6, the lexical pipeline achieved a precision of 0.9015, recall of 0.7933, F1-score of 0.8440, and false-positive rate of 0.0867.
\clearpage

\tableofcontents
\clearpage
\phantomsection
\addcontentsline{toc}{chapter}{Abbildungsverzeichnis}
\listoffigures
\clearpage
\phantomsection
\addcontentsline{toc}{chapter}{Tabellenverzeichnis}
\listoftables
\clearpage
\phantomsection
\addcontentsline{toc}{chapter}{Quell- und Pseudocodeverzeichnis}
\lstlistoflistings
\clearpage

\chapter*{Abkürzungsverzeichnis, Fachbegriffe und Einheiten}
\addcontentsline{toc}{chapter}{Abkürzungsverzeichnis, Fachbegriffe und Einheiten}
\begin{longtable}{|L{4cm}|L{10cm}|}
\hline
\textbf{Abkürzung / Begriff} & \textbf{Bedeutung}\\
\hline
API & Application Programming Interface\\
\hline
Base64 & Kodierungsverfahren zur Darstellung binärer Daten als ASCII-Zeichen\\
\hline
C2 & Command and Control; in dieser Arbeit nur als ungefährliches, simuliertes Muster betrachtet\\
\hline
CLI & Command Line Interface\\
\hline
EDR & Endpoint Detection and Response\\
\hline
Eventlog & Windows-Ereignisprotokoll\\
\hline
F1 & Harmonisches Mittel aus Precision und Recall\\
\hline
FPR & False-Positive-Rate\\
\hline
HIDS & Host-based Intrusion Detection System\\
\hline
IEX & PowerShell-Alias für \texttt{Invoke-Expression}\\
\hline
JSON & JavaScript Object Notation\\
\hline
LOLBins & Living-off-the-Land Binaries; legitime Systemwerkzeuge, die missbraucht werden können\\
\hline
PowerShell & Windows-Shell und Skriptumgebung zur Administration und Automatisierung\\
\hline
Recall & Anteil korrekt erkannter positiver Fälle an allen tatsächlich positiven Fällen\\
\hline
ScriptBlock & Von PowerShell interpretierter Skriptabschnitt\\
\hline
SIEM & Security Information and Event Management\\
\hline
Wazuh & Open-Source-Plattform für Sicherheitsüberwachung und hostbasierte Erkennung\\
\hline
XML & Extensible Markup Language\\
\hline
\end{longtable}
\clearpage

%--------------------------------------------------------------------
% Hauptteil
%--------------------------------------------------------------------
\pagenumbering{arabic}

\chapter{Einleitung}

\section{Motivation und Hintergrund}
PowerShell ist unter Windows ein zentrales Werkzeug für Administration, Automatisierung und Systemverwaltung und wird im MITRE-ATT\&CK-Framework als relevante Technik für die Ausführung von Befehlen und Skripten beschrieben \cite{mitre-powershell}. Administratorinnen und Administratoren nutzen PowerShell, um Systeme zu konfigurieren, Ereignisse auszuwerten, Dienste zu steuern oder wiederkehrende Aufgaben zu automatisieren. Dieselben Eigenschaften machen PowerShell jedoch auch für Angreifende interessant. Skripte können direkt über legitime Windows-Komponenten gestartet, dynamisch zusammengesetzt, kodiert oder im Arbeitsspeicher ausgeführt werden. Für Sicherheitssysteme entsteht dadurch die Herausforderung, legitime Administration von missbräuchlicher Skriptausführung zu unterscheiden.

Ein besonders relevantes Problem ist die Obfuskierung von PowerShell-Code. Bereits einfache Techniken wie Base64-Kodierung, String-Konkatenation, Backticks oder wechselnde Groß- und Kleinschreibung können einfache Signaturen umgehen. Gleichzeitig stellt Windows mit dem ScriptBlock-Logging eine wichtige Datenquelle bereit. Wenn ScriptBlock-Logging aktiviert ist, können Inhalte von PowerShell-Ausführungen im Windows-Eventlog protokolliert und anschließend von einem Host Intrusion Detection System wie Wazuh ausgewertet werden \cite{elastic-powershell-winlogbeat,splunk-scriptblock}.

Die Arbeit setzt an dieser Schnittstelle an. Sie untersucht, ob Wazuh auf Basis von ScriptBlock-Logging ausreichend Hinweise auf obfuskierte PowerShell-Aktivität liefert und ob eine zusätzliche, einfache lexikalische Analyse die Erkennung verbessern kann. Der Fokus liegt dabei nicht auf einem vollständigen Produktivsystem, sondern auf einem nachvollziehbaren und reproduzierbaren Laboraufbau.

\section{Problemstellung}
Regelbasierte Erkennung ist transparent und gut erklärbar. Eine Regel kann beispielsweise auf bestimmte Tokens, Parameter oder Kombinationen von PowerShell-Befehlen reagieren. Der Nachteil besteht darin, dass Regeln häufig konkrete Muster erwarten. Wird ein verdächtiger Befehl durch String-Konkatenation, Aliase oder Kodierung verändert, kann die Regel versagen. Gleichzeitig sind zu breite Regeln problematisch, weil sie legitime administrative Tätigkeiten fälschlicherweise als Angriff markieren können.

Die zentrale Problemstellung lautet daher: Eine gute Erkennung muss einerseits robust gegenüber einfachen Obfuskierungen sein und andererseits die FPR bei legitimen Skripten gering halten. Diese Arbeit untersucht diese Abwägung anhand von drei Vergleichsgruppen: Wazuh allein, lexikalische Analyse allein und eine Kombination beider Ansätze.

\section{Zielsetzung der Arbeit}
Ziel der Arbeit ist die Konzeption, Implementierung und Auswertung eines hostbasierten Erkennungs- und Evaluierungssystems für obfuskierte PowerShell-Aktivität unter Windows 11. Dazu werden ScriptBlock-Logs gesammelt, durch Wazuh ausgewertet und zusätzlich mit einer Python-Pipeline analysiert. Die Python-Pipeline extrahiert einfache lexikalische Merkmale aus PowerShell-ScriptBlocks und vergleicht deren Erkennungsleistung mit der regelbasierten Wazuh-Erkennung.

Die Arbeit verfolgt folgende konkrete Ziele:
\begin{enumerate}[label=Z\arabic*]
    \item Aufbau einer kontrollierten Windows-11-Testumgebung mit aktiviertem ScriptBlock-Logging und Wazuh-Anbindung.
    \item Erstellung eines kuratierten Datensatzes aus legitimen administrativen PowerShell-Skripten, simulierten Angriffsmustern und obfuskierten Varianten.
    \item Entwicklung einer regelbasierten Wazuh-Baseline auf Basis vorhandener und eigener Regeln.
    \item Implementierung einer Python-Pipeline zur Extraktion lexikalischer Merkmale.
    \item Vergleich der Erkennungsansätze anhand von Precision, Recall, F1-Score und FPR.
\end{enumerate}

\section{Forschungsfragen}
Die zentrale Forschungsfrage lautet:

\begin{quote}
Wie gut erkennt Wazuh obfuskierte PowerShell-Aktivität über ScriptBlock-Logging, und kann eine einfache lexikalische Analyse die Erkennung verbessern?
\end{quote}

Daraus ergeben sich folgende Teilfragen:
\begin{enumerate}[label=F\arabic*]
    \item Welche obfuskierten PowerShell-Aktivitäten werden durch Standard- und eigene Wazuh-Regeln erkannt?
    \item Welche lexikalischen Merkmale unterscheiden legitime administrative ScriptBlocks von simulierten maliziösen ScriptBlocks?
    \item Verbessert die Kombination aus Wazuh-Regeln und lexikalischen Merkmalen Precision, Recall oder F1-Score gegenüber den Einzelverfahren?
    \item Wie hoch ist die FPR bei legitimen administrativen PowerShell-Skripten?
    \item Welche Obfuskierungstechniken führen besonders häufig zu False Negatives oder zu auffälligen lexikalischen Profilen?
\end{enumerate}

\section{Abgrenzung und Sicherheitsrahmen}
Die Experimente werden ausschließlich in einer lokalen, kontrollierten Laborumgebung durchgeführt. C2-ähnliche PowerShell-Muster werden in die Tests aufgenommen, da sie für die Erkennung maliziöser ScriptBlocks relevant sind. Dabei werden jedoch keine echten externen Command-and-Control-Verbindungen aufgebaut und keine funktionsfähigen Malware-Payloads nachgeladen oder ausgeführt. Netzwerknahe Angriffsmuster werden stattdessen kontrolliert simuliert, etwa durch Platzhalter-URLs, lokale Testserver, nicht erreichbare Testadressen oder nicht ausführende Strings.

Damit gilt für diese Arbeit: C2-Indikatoren ja, echte C2-Infrastruktur nein. Die Arbeit bewertet außerdem keine vollständige Endpoint-Detection-and-Response-Lösung und keine produktive Malware-Erkennung. Sie konzentriert sich auf ScriptBlock-Logging, Wazuh-Regeln, einfache lexikalische Merkmale und eine reproduzierbare experimentelle Auswertung.

\section{Vorgehensweise}
Die Vorgehensweise orientiert sich an experimentellen Bachelorarbeiten mit klarer Trennung zwischen Grundlagen, Design, Implementierung und Evaluation. Zunächst werden die theoretischen Grundlagen und der Stand der Technik beschrieben. Anschließend wird die Testumgebung entworfen und implementiert. Darauf folgt die Erstellung eines Datensatzes mit legitimen und simulierten maliziösen PowerShell-Samples. Die Experimente vergleichen drei Erkennungsvarianten und werten die Ergebnisse anhand einheitlicher Metriken aus.

\section{Aufbau der Arbeit}
Kapitel 2 erläutert die Grundlagen zu PowerShell, ScriptBlock-Logging, Obfuskierung und Wazuh. Kapitel 3 beschreibt den Stand der Technik und grenzt regelbasierte, anomaliebasierte und lexikalische Erkennung voneinander ab. Kapitel 4 formuliert Anforderungen und das Design der Testumgebung. Kapitel 5 beschreibt die Implementierung der Laborumgebung, Wazuh-Regeln und Python-Pipeline. Kapitel 6 erläutert die Evaluation und die Ergebnisdarstellung. Kapitel 7 diskutiert Grenzen, Validität und praktische Empfehlungen. Kapitel 8 fasst die Arbeit zusammen und gibt einen Ausblick. Der Anhang enthält Datensatzschema, Regelvorlagen, Pseudocode und Protokolle, die die Nachvollziehbarkeit der Experimente unterstützen.

\chapter{Grundlagen}

\section{PowerShell als Administrations- und Automatisierungsumgebung}
PowerShell ist eine Shell- und Skriptumgebung, die besonders unter Windows tief in Systemverwaltung und Automatisierung eingebunden ist. Sie bietet Zugriff auf Dateien, Prozesse, Dienste, Registry, Netzwerkfunktionen, Windows Management Instrumentation und .NET-Klassen. Dadurch lassen sich viele administrative Aufgaben effizient automatisieren.

Aus Sicht der IT-Sicherheit ist diese Leistungsfähigkeit ambivalent. PowerShell ist auf vielen Windows-Systemen vorhanden und muss nicht nachinstalliert werden. Angreifende können daher legitime Systemfunktionen missbrauchen, ohne auffällige zusätzliche Binärdateien abzulegen. Dieses Prinzip wird häufig als Living-off-the-Land bezeichnet. Für die Erkennung bedeutet dies, dass die bloße Nutzung von PowerShell nicht automatisch verdächtig ist. Entscheidend sind Kontext, Befehlskette, Parameter, Tokens, Kodierung und Ausführungsumgebung.

\section{ScriptBlock-Logging}
ScriptBlock-Logging ist eine Windows-PowerShell-Protokollierungsfunktion, mit der ausgeführte ScriptBlocks im Windows-Ereignisprotokoll gespeichert werden können \cite{microsoft-logging-windows,microsoft-logging}. Für diese Arbeit ist diese Telemetriequelle zentral, weil sie deutlich mehr semantische Informationen enthält als reine Prozessstartdaten. Während ein Prozessstartlog häufig nur zeigt, dass \texttt{powershell.exe} oder \texttt{pwsh.exe} gestartet wurde, enthält ein ScriptBlock-Log den interpretierten Skriptinhalt oder relevante Teile davon.

Die Qualität der späteren Erkennung hängt stark davon ab, ob ScriptBlock-Logging korrekt aktiviert ist und ob die Ereignisse vollständig an Wazuh weitergeleitet werden. Daher wird die Aktivierung im Methodik- und Implementierungsteil genau dokumentiert.

\screenshotplaceholder{powershell_scriptblock_event_4104.png}{Beispiel eines PowerShell-ScriptBlock-Ereignisses im Windows Event Viewer}{fig:scriptblock-event-4104}

\section{PowerShell-Obfuskierung}
Obfuskierung bezeichnet die absichtliche Veränderung eines Skripts, sodass dessen Zweck für Menschen oder einfache Erkennungsregeln schlechter erkennbar wird, während die Funktionalität erhalten bleibt. In dieser Arbeit werden einfache und reproduzierbare Techniken betrachtet, weil diese für eine Bachelorarbeit transparent analysierbar sind.

\begin{table}[H]
\centering
\caption{Betrachtete PowerShell-Obfuskierungstechniken}
\begin{tabularx}{\textwidth}{|L{5cm}|Y|}
\hline
\textbf{Technik} & \textbf{Beschreibung und Relevanz}\\
\hline
Base64-Encoding & Kodiert Befehle und tritt häufig zusammen mit \texttt{-EncodedCommand} oder Dekodierfunktionen auf.\\
\hline
String-Konkatenation & Verdächtige Tokens werden aus Teilstrings zusammengesetzt, zum Beispiel aus getrennten Wortbestandteilen.\\
\hline
Variablenverschleierung & Aussagekräftige Variablennamen werden durch kurze, zufällige oder irreführende Namen ersetzt.\\
\hline
Aliase & Ausgeschriebene Cmdlets werden durch kurze Aliasformen ersetzt.\\
\hline
Backticks & Escape-Zeichen können Tokens aufbrechen und einfache Suchmuster erschweren.\\
\hline
Whitespace-Veränderung & Leerzeichen, Tabs und Zeilenumbrüche werden ungewöhnlich platziert.\\
\hline
Gemischte Groß-/Kleinschreibung & Schreibweisen werden variiert, um case-sensitive Regeln zu umgehen.\\
\hline
Command Reconstruction & Befehle werden erst zur Laufzeit aus Zeichenketten rekonstruiert.\\
\hline
\end{tabularx}
\end{table}

\section{Hostbasierte Intrusion Detection}
Hostbasierte Intrusion Detection Systeme überwachen Ereignisse auf einzelnen Endsystemen. NIST beschreibt Host-basierte Verfahren als eine zentrale Kategorie von Intrusion-Detection- und Intrusion-Prevention-Systemen \cite{nist-idps}. Dazu zählen Systemlogs, Prozessinformationen, Konfigurationsänderungen, Dateiänderungen und sicherheitsrelevante Ereignisse. Im Gegensatz zu rein netzwerkbasierten Verfahren können hostbasierte Ansätze Aktivitäten erkennen, die lokal stattfinden oder durch verschlüsselte Kommunikation von außen nicht sichtbar wären.

Für PowerShell-Erkennung sind hostbasierte Daten besonders wichtig, weil viele verdächtige Aktivitäten direkt auf dem Endpunkt ausgeführt werden. ScriptBlock-Logging stellt dabei eine inhaltsnahe Datenquelle dar, die sich mit Wazuh-Regeln oder zusätzlichen Analyseverfahren auswerten lässt.

\section{Wazuh}
Wazuh ist eine Open-Source-Plattform zur Sicherheitsüberwachung. Die Wazuh-Dokumentation beschreibt unter anderem die Sammlung von Windows-Eventlogs sowie die XML-basierte Regelsyntax für eigene Regeln \cite{wazuh-log-collection,wazuh-rules,wazuh-custom-rules}. Sie kombiniert Agenten auf Endsystemen mit einem zentralen Manager und einem Dashboard. In dieser Arbeit wird Wazuh eingesetzt, um Windows-Ereignisse des Testsystems zu sammeln, PowerShell-relevante Ereignisse zu analysieren und daraus Alerts zu erzeugen.

Wazuh bietet eine regelbasierte Erkennung. Diese ist gut nachvollziehbar, weil ausgelöste Alerts auf konkrete Regeln zurückgeführt werden können. Gleichzeitig ist sie anfällig für Umgehungen, wenn ein Angriffsmuster nicht mehr zum erwarteten Regelmuster passt. Genau an dieser Stelle setzt die ergänzende lexikalische Analyse an.

\section{Lexikalische Merkmale}
Lexikalische Merkmale beschreiben Eigenschaften eines Textes, ohne dessen vollständige Semantik zu verstehen. Bei PowerShell-ScriptBlocks können solche Merkmale Hinweise auf Obfuskierung oder ungewöhnliche Ausführungsmuster liefern. Beispiele sind Länge, Sonderzeichendichte, Token-Häufigkeiten, Entropie, Anzahl von Backticks, Anzahl von Konkatenationen oder das Auftreten verdächtiger Schlüsselwörter.

Der Vorteil lexikalischer Merkmale liegt in ihrer einfachen Berechenbarkeit und Erklärbarkeit. Der Nachteil besteht darin, dass sie keine Garantie für Malizität liefern. Ein legitimes Administrationsskript kann lang oder komplex sein, während ein gefährlicher Befehl sehr kurz sein kann. Daher wird in dieser Arbeit nicht nur die Erkennungsrate, sondern auch die FPR betrachtet.

\chapter{Stand der Technik und verwandte Ansätze}

\section{Regelbasierte Erkennung}
Regelbasierte Erkennung sucht nach bekannten Mustern. Diese können einfache Zeichenketten, reguläre Ausdrücke, Kombinationen von Feldern oder komplexere Bedingungen sein. Der Vorteil liegt in der Nachvollziehbarkeit: Wenn eine Regel auslöst, ist klar, welches Muster erkannt wurde. Dies erleichtert Dokumentation, Fehlersuche und Anpassung.

Der Nachteil besteht darin, dass Regeln bekannte oder erwartete Muster voraussetzen. Wenn ein Befehl kodiert, rekonstruiert oder syntaktisch verändert wird, kann eine Regel nicht mehr greifen. Für obfuskierte PowerShell-Aktivität ist diese Grenze besonders relevant, weil schon einfache Textveränderungen ausreichen können.

\section{Anomaliebasierte Erkennung}
Anomaliebasierte Verfahren definieren zunächst ein Normalverhalten und markieren Abweichungen. Dadurch können auch unbekannte Muster auffallen. Gleichzeitig sind solche Ansätze anfälliger für False Positives, wenn legitime administrative Tätigkeit vom gelernten Normalverhalten abweicht. Für eine Bachelorarbeit mit begrenztem Datensatz kann ein vollständig anomaliebasiertes Modell zu aufwendig sein.

Diese Arbeit nutzt deshalb keine komplexe Anomalieerkennung, sondern eine einfache lexikalische Analyse als transparenten Zwischenweg. Sie ist erklärbarer als ein komplexes Machine-Learning-Modell und kann trotzdem Hinweise liefern, die über konkrete Signaturen hinausgehen.

\section{Vergleich mit IDS-orientierten Referenzarbeiten}
Experimentelle IDS-Arbeiten sind häufig so aufgebaut, dass zunächst die theoretischen Grundlagen beschrieben werden, anschließend die Testumgebung und dann konkrete Angriffsszenarien beziehungsweise Testfälle. Danach folgt eine Auswertung pro Testfall, Methode und Metrik. Diese Struktur wird für die vorliegende Arbeit übernommen: Grundlagen, Methodik, Testumgebung, Implementierung, Evaluation, Diskussion und Ausblick.

Besonders wichtig ist die saubere Trennung zwischen Testdurchführung und Interpretation. Messwerte gehören in das Evaluationskapitel, während die Bedeutung der Ergebnisse, Grenzen und Empfehlungen in der Diskussion behandelt werden.

\section{Einordnung der eigenen Arbeit}
Die Arbeit ist weder eine reine Literaturarbeit noch eine reine Implementierungsarbeit. Sie kombiniert Laboraufbau, Datensatzkuration, Regelentwicklung, Python-Implementierung und quantitative Auswertung. Die wissenschaftliche Leistung entsteht durch die reproduzierbare Gegenüberstellung der drei Erkennungsgruppen und durch die Analyse, welche Obfuskierungsarten welche Erkennung beeinflussen.

\chapter{Anforderungen und Design}

\section{Funktionale Anforderungen}
Die folgenden funktionalen Anforderungen leiten sich aus der Zielsetzung ab:
\begin{enumerate}[label=FA\arabic*]
    \item Das Windows-11-Testsystem muss PowerShell-ScriptBlocks protokollieren.
    \item Der Wazuh-Agent muss relevante Windows-Ereignisse an den Wazuh-Manager übertragen.
    \item Wazuh muss Standardregeln und eigene Regeln auf PowerShell-Ereignisse anwenden.
    \item Die Testfälle müssen reproduzierbar ausgeführt und eindeutig gelabelt werden können.
    \item Die Python-Pipeline muss ScriptBlock-Logs einlesen und Features extrahieren können.
    \item Die Auswertung muss Metriken für Wazuh, lexikalische Analyse und Kombination berechnen.
\end{enumerate}

\section{Nicht-funktionale Anforderungen}
Neben der Funktionalität sind folgende nicht-funktionale Anforderungen relevant:
\begin{enumerate}[label=NFA\arabic*]
    \item \textbf{Reproduzierbarkeit:} Konfiguration, Versionen und Testfälle müssen dokumentiert werden.
    \item \textbf{Sicherheit:} Es dürfen keine echten externen C2-Verbindungen und keine funktionsfähigen Malware-Payloads ausgeführt werden.
    \item \textbf{Nachvollziehbarkeit:} Regeln, Features und Schwellenwerte müssen transparent beschrieben werden.
    \item \textbf{Trennung der Daten:} Rohdaten, bereinigte Daten, Features und Ergebnisse sollen getrennt gespeichert werden.
    \item \textbf{Erweiterbarkeit:} Neue Obfuskierungstechniken oder Features sollen später ergänzt werden können.
\end{enumerate}

\section{Anforderungsstatus}
Die Anforderungen werden zusätzlich als Statusübersicht dokumentiert. Dadurch ist sichtbar, welche Bestandteile umgesetzt wurden und welche Ergebnisse aus dem finalen Wazuh-Alert-Export stammen.

\begin{table}[H]
\centering
\caption{Anforderungen mit Umsetzungsstatus}
\small
\begin{tabularx}{\textwidth}{|C{2cm}|Y|C{3cm}|}
\hline
\textbf{ID} & \textbf{Beschreibung} & \textbf{Status}\\
\hline
FA1 & PowerShell-ScriptBlocks werden auf dem Windows-Testsystem protokolliert. & umgesetzt\\
\hline
FA2 & Der Wazuh-Agent ist installiert und überträgt relevante Windows-Ereignisse an den Wazuh-Manager. & umgesetzt\\
\hline
FA3 & Standardmäßige Wazuh-Regeln werden als Baseline gegen den Datensatz getestet. & umgesetzt\\
\hline
FA4 & Jeder Testfall ist eindeutig gelabelt und über eine \texttt{WAZUH\_SAMPLE}-ID zuordenbar. & umgesetzt\\
\hline
FA5 & Die Python-Pipeline extrahiert lexikalische Merkmale aus den ScriptBlocks. & umgesetzt\\
\hline
FA6 & Die Pipeline berechnet Precision, Recall, F1-Score und FPR für Wazuh, lexikalische Analyse und Kombination. & umgesetzt\\
\hline
NFA1 & Datensatz, Skripte, Schwellenwerte und Ergebnisdateien sind reproduzierbar dokumentiert. & umgesetzt\\
\hline
NFA2 & Es werden keine realen Malware-Payloads oder externen C2-Verbindungen ausgeführt. & umgesetzt\\
\hline
NFA3 & Der finale Vergleich darf erst nach Import echter Wazuh-Alerts als abgeschlossen gelten. & umgesetzt\\
\hline
\end{tabularx}
\end{table}

\section{Lab-Architektur}
Die Testumgebung besteht aus einem überwachten Windows-11-System und einer Wazuh-Instanz. Die Wazuh-Instanz wird in Docker oder in einer separaten virtuellen Maschine betrieben. Die Python-Auswertung findet auf dem Hostsystem oder auf einer separaten Analyse-VM statt.

\begin{figure}[H]
\centering
\includegraphics[width=0.95\textwidth]{figures/lab_architecture_activity.png}
\caption{Konzeptionelle Lab-Architektur der Arbeit}
\label{fig:lab-architecture}
\end{figure}

\begin{table}[H]
\centering
\caption{Komponenten der Testumgebung}
\begin{tabularx}{\textwidth}{|L{4cm}|Y|}
\hline
\textbf{Komponente} & \textbf{Aufgabe}\\
\hline
Windows 11 VM/Host & Ausführung legitimer und simulierter maliziöser PowerShell-Skripte.\\
\hline
PowerShell ScriptBlock-Logging & Erfassung der ausgeführten ScriptBlocks im Windows-Eventlog.\\
\hline
Wazuh Agent & Weiterleitung relevanter Windows-Ereignisse an den Wazuh-Manager.\\
\hline
Wazuh Manager & Regelbasierte Analyse und Alert-Erzeugung.\\
\hline
Wazuh Dashboard & Sichtung, Filterung und Export der Alerts.\\
\hline
Python-Pipeline & Feature-Extraktion, Klassifikation und Metrikberechnung.\\
\hline
\end{tabularx}
\end{table}

\begin{table}[H]
\centering
\caption{Technische Testumgebung und verwendete Versionen}
\small
\begin{tabularx}{\textwidth}{|L{4cm}|Y|}
\hline
\textbf{Komponente} & \textbf{Version beziehungsweise Ausprägung}\\
\hline
Windows-Testsystem & Überwachtes Windows-System \texttt{MSI}; lokale Systemabfrage der Bearbeitungsumgebung: Windows 10 Home 25H2, Build 26200.8737.\\
\hline
PowerShell & Windows PowerShell 5.1.26100.8737.\\
\hline
Wazuh-Agent & Wazuh Agent 4.14.4 rc2; lokaler Dienst \texttt{WazuhSvc} läuft.\\
\hline
Wazuh Manager und Dashboard & Wazuh Server v4.14.5 rc1 in WSL Kali auf derselben physischen Maschine; Alert-Datei \texttt{/var/ossec/logs/alerts/alerts.json}.\\
\hline
Docker/VM-Aufbau & Gemessener Lauf ohne Docker-Manager: Windows-Agent auf dem Hostsystem, Wazuh-Manager in WSL Kali.\\
\hline
Python & Python 3.12.13 für die Ausführung der Pipeline.\\
\hline
CPU und Arbeitsspeicher & 12th Gen Intel(R) Core(TM) i7-1280P, ca. 16 GB RAM laut lokaler Systemabfrage.\\
\hline
Datensatz & \codebreakthree{powershell\_scriptblock}{\_samples\_final}{\_thesis\_v4.json}, 300 Samples mit 150 benignen und 150 simuliert maliziösen Fällen.\\
\hline
\end{tabularx}
\end{table}

\screenshotplaceholder{lab_environment_versions.png}{Dokumentation der finalen Laborversionen und Systeminformationen}{fig:lab-environment-versions}

\section{Datensatzdesign}
Der Datensatz bildet die Grundlage der Evaluation und liegt als JSON-Datei in der Version \texttt{powershell\_scriptblock\_samples\_final\_thesis\_v4.json} vor. Jeder Eintrag enthält eine eindeutige ID, ein Label, eine Szenariobeschreibung, die taktische Einordnung, die zugehörige MITRE-ATT\&CK-Technik, den ScriptBlock, den Obfuskierungstyp, Sicherheitsnotizen sowie Felder für spätere Wazuh- und Pipeline-Ergebnisse. Jeder ausgeführte ScriptBlock enthält zusätzlich eine Markierung der Form \texttt{WAZUH\_SAMPLE <ID> <Szenario>}. Dadurch können exportierte Wazuh-Alerts später wieder eindeutig einem Datensatz-Eintrag zugeordnet werden.

\begin{table}[H]
\centering
\caption{Datensatzschema der Version 4}
\small
\begin{tabularx}{\textwidth}{|L{4cm}|Y|}
\hline
\textbf{Feld} & \textbf{Bedeutung}\\
\hline
\texttt{id} & Eindeutige Sample-ID, zum Beispiel \texttt{BEN-001}, \texttt{SIM-001} oder \texttt{OBF-001}.\\
\hline
\texttt{label} & Klassenlabel \texttt{benign} oder \texttt{simulated\_malicious}.\\
\hline
\texttt{dataset\_group} & Gruppierung in normale benigne, schwere benigne, klare Indikator- oder obfuskierte Indikator-Samples.\\
\hline
\texttt{scenario} & Kurzbeschreibung des betrachteten Testfalls.\\
\hline
\codebreak{mitre\_attack}{\_id} & MITRE-ATT\&CK-Technik, falls zutreffend, sonst \texttt{N/A}.\\
\hline
\codebreak{obfuscation}{\_type} & Verwendete Obfuskierungsart oder \texttt{none}.\\
\hline
\texttt{scriptblock} & PowerShell-ScriptBlock, der in der Laborumgebung ausgeführt beziehungsweise protokolliert wird.\\
\hline
\codebreak{safety}{\_class} & Sicherheitsklassifikation, etwa \texttt{benign\_safe\_execution} oder \codebreakthree{simulated}{\_malicious}{\_non\_executable}.\\
\hline
\texttt{features} & Von der Python-Pipeline berechnete lexikalische Merkmale.\\
\hline
\codebreak{wazuh\_alert}{\_observed} & Ergebnisfeld für die spätere Wazuh-Auswertung.\\
\hline
\codebreak{lexical}{\_score} & Von der Python-Pipeline berechneter Score.\\
\hline
\codebreak{lexical}{\_prediction} & Entscheidung der lexikalischen Pipeline.\\
\hline
\end{tabularx}
\end{table}

\section{Klassen und Testfälle}
Der finale Datensatz umfasst 300 Samples mit ausgeglichener Klassenverteilung. Die positive Klasse besteht ausschließlich aus simuliert maliziösen, nicht-ausführenden Indikator-Samples. Dadurch können verdächtige PowerShell-Muster und Obfuskierungen untersucht werden, ohne reale Malware, echte C2-Verbindungen oder produktive Angriffsinfrastruktur zu verwenden.

\begin{table}[H]
\centering
\caption{Klassen und Gruppen des finalen Datensatzes}
\small
\begin{tabularx}{\textwidth}{|L{3cm}|L{2.5cm}|C{1.5cm}|Y|}
\hline
\textbf{Gruppe} & \textbf{Label} & \textbf{Anzahl} & \textbf{Beschreibung}\\
\hline
Normale benigne Samples & benign & 100 & Typische administrative PowerShell-Befehle, etwa Prozess-, Dienst-, Eventlog- und Systeminventarisierung.\\
\hline
Schwere benigne Samples & benign & 50 & Legitime Skripte mit höherer lexikalischer Auffälligkeit, etwa längere Strings, interne URLs, Formatierung, JSON/CSV/XML-Verarbeitung oder sichere Base64-Konfigurationstexte.\\
\hline
Direkte Indikator-Samples & \codebreak{simulated}{\_malicious} & 30 & Klar erkennbare, aber nicht-ausführende Indikator-ScriptBlocks mit verdächtigen PowerShell-, LOLBin- oder C2-ähnlichen Begriffen.\\
\hline
Obfuskierte Indikator-Samples & \codebreak{simulated}{\_malicious} & 120 & Varianten der simuliert maliziösen Indikatoren mit Base64, String-Konkatenation, Backticks, gemischter Schreibweise, Alias-/Befehlsrekonstruktion und kombinierter Obfuskation.\\
\hline
\end{tabularx}
\end{table}

\begin{table}[H]
\centering
\caption{Obfuskierungsarten im finalen Datensatz}
\small
\begin{tabularx}{\textwidth}{|L{4cm}|C{1.5cm}|Y|}
\hline
\textbf{Obfuskierungsart} & \textbf{Anzahl} & \textbf{Ziel der Testfälle}\\
\hline
\codebreak{base64\_literal}{\_decode\_safe} & 20 & Prüfung, ob Kodierungsindikatoren und hohe Entropie erkannt werden.\\
\hline
\codebreak{string}{\_concatenation} & 20 & Prüfung, ob aufgeteilte Tokens und Konkatenationsoperatoren auffallen.\\
\hline
\codebreak{backtick}{\_literal} & 20 & Prüfung, ob PowerShell-spezifische Token-Unterbrechung erkannt wird.\\
\hline
\codebreak{mixed}{\_case} & 20 & Prüfung der Robustheit gegenüber wechselnder Groß- und Kleinschreibung.\\
\hline
\codebreak{alias\_command}{\_reconstruction} & 20 & Prüfung von Aliasnutzung und rekonstruierten Befehlen.\\
\hline
\codebreak{combined}{\_obfuscation} & 20 & Prüfung mehrschichtiger Obfuskierung aus mehreren einfachen Techniken.\\
\hline
\end{tabularx}
\end{table}

\section{Sicheres Design der C2-ähnlichen Muster}
Da C2-Muster für PowerShell-Erkennung relevant sind, werden sie als harmlose Indikatoren in die Experimente aufgenommen. Es werden jedoch keine realen externen Ziele kontaktiert. Folgende Varianten sind zulässig:
\begin{itemize}
    \item Platzhalter-URLs wie \texttt{http://example.local/test.ps1},
    \item lokale Testadressen wie \texttt{127.0.0.1},
    \item nicht ausführende Strings, die nur geloggt werden,
    \item lokale Dateien mit harmlosen Testinhalten,
    \item Testserver ohne Schadfunktion, falls Netzwerkverhalten lokal beobachtet werden soll.
\end{itemize}

\begin{lstlisting}[language={},caption={Harmloses Beispiel eines C2-ähnlichen Strings ohne Ausführung}]
$indicator = 'IEX (New-Object Net.WebClient).DownloadString("http://example.local/payload.ps1")'
Write-Output $indicator
\end{lstlisting}

\section{Auswertungsdesign}
Die Erkennung wird in drei Gruppen verglichen:
\begin{enumerate}
    \item \textbf{Nur Wazuh:} Es zählt ausschließlich, ob Wazuh einen Alert erzeugt.
    \item \textbf{Nur lexikalische Analyse:} Es zählt ausschließlich die Entscheidung der Python-Pipeline.
    \item \textbf{Kombination:} Ein Eintrag gilt als erkannt, wenn Wazuh oder die lexikalische Analyse anschlägt.
\end{enumerate}

Diese Struktur ermöglicht einen direkten Vergleich der Erkennungsleistung und zeigt, ob die lexikalische Analyse einen Zusatznutzen gegenüber Wazuh allein liefert.

\screenshotplaceholder{pipeline_flow_diagram.png}{Grafische Darstellung des vollständigen Experimentablaufs}{fig:pipeline-flow-placeholder}

\chapter{Implementierung}

\section{Einrichtung der Testumgebung}
Die Implementierung beginnt mit dem Aufbau der Laborumgebung. Das Windows-11-System wird als überwachte Maschine eingerichtet. Danach wird ScriptBlock-Logging aktiviert und überprüft, ob PowerShell-Ausführungen im Windows-Eventlog erscheinen. Anschließend wird der Wazuh-Agent installiert, mit dem Wazuh-Manager verbunden und so konfiguriert, dass relevante Windows-Eventlogs weitergeleitet werden.

Die Wazuh-Instanz wird entweder über Docker oder in einer separaten VM betrieben. Wichtig ist, dass die Umgebung isoliert bleibt und keine Testaktivitäten unbeabsichtigt in externe Netze gelangen.

\screenshotplaceholder{wazuh_agent_status.png}{Status des Windows-Agents im Wazuh-Dashboard}{fig:wazuh-agent-status}

\screenshotplaceholder{wazuh_dashboard.png}{Wazuh-Dashboard mit gefilterten PowerShell-Alerts}{fig:wazuh-dashboard}

\section{Aktivierung von ScriptBlock-Logging}
Die Aktivierung kann über Gruppenrichtlinien oder Registry-Einstellungen erfolgen. In dieser Arbeit wird der konkrete Aktivierungsweg dokumentiert. Nach der Aktivierung werden Testbefehle ausgeführt und die erzeugten Ereignisse kontrolliert.

\screenshotplaceholder{windows_scriptblock_gpo.png}{Aktivierung von PowerShell ScriptBlock-Logging über Gruppenrichtlinien oder Registry}{fig:windows-scriptblock-gpo}

\screenshotplaceholder{windows_event_viewer_scriptblock.png}{Kontrolle eines erzeugten ScriptBlock-Events im Windows Event Viewer}{fig:windows-event-viewer-scriptblock}

\screenshotplaceholder{windows_event_viewer_wazuh_sample_marker.png}{Windows Event Viewer mit sichtbarer \texttt{WAZUH\_SAMPLE}-Markierung eines Testfalls}{fig:windows-event-viewer-wazuh-sample}

\begin{table}[H]
\centering
\caption{Checkliste zur Logging-Verifikation}
\begin{tabularx}{\textwidth}{|L{4cm}|Y|}
\hline
\textbf{Prüfschritt} & \textbf{Erwartetes Ergebnis}\\
\hline
ScriptBlock-Logging aktiviert & PowerShell-ScriptBlocks erscheinen im Eventlog.\\
\hline
Testskript ausgeführt & Ereignis enthält den erwarteten ScriptBlock-Inhalt.\\
\hline
Wazuh-Agent verbunden & Agent erscheint im Wazuh-Dashboard als aktiv.\\
\hline
Eventweiterleitung aktiv & PowerShell-Ereignisse sind in Wazuh auffindbar.\\
\hline
Alert-Erzeugung geprüft & Mindestens eine Testregel erzeugt einen nachvollziehbaren Alert.\\
\hline
\end{tabularx}
\end{table}

\section{Erstellung legitimer administrativer Skripte}
Legitime Skripte sollen typische Administrationsaufgaben abbilden. Dazu gehören Systeminformationen, Prozess- und Dienstlisten, Eventlog-Abfragen, Benutzer- und Gruppeninformationen sowie einfache Dateisystemoperationen. Diese Skripte sind wichtig, um False Positives zu messen.

\begin{lstlisting}[language={},caption={Beispielkategorien legitimer administrativer PowerShell-Skripte}]
Get-Process
Get-Service
Get-EventLog -LogName System -Newest 20
Get-LocalUser
Get-ChildItem -Path C:\Windows -Recurse -ErrorAction SilentlyContinue
\end{lstlisting}

\section{Erstellung simulierter maliziöser Muster}
Die simulierten Muster enthalten typische verdächtige Tokens, führen aber keine Schadfunktion aus. Dadurch kann untersucht werden, ob Wazuh und die Python-Pipeline bekannte Indikatoren erkennen, ohne eine reale Angriffsaktivität auszuführen.

Beispiele für untersuchte Indikatoren sind:
\begin{itemize}
    \item \texttt{Invoke-Expression} und \texttt{IEX},
    \item \texttt{DownloadString}, \texttt{WebClient} und \texttt{Invoke-WebRequest},
    \item \texttt{FromBase64String},
    \item \texttt{-EncodedCommand},
    \item dynamische Zusammensetzung von Befehlen.
\end{itemize}

\section{Wazuh-Regeln}
Für die Wazuh-Baseline werden zunächst Standardregeln untersucht. Danach werden eigene Regeln ergänzt. Diese eigenen Regeln erkennen typische PowerShell-Muster und dokumentieren, welche Indikatoren jeweils abgedeckt sind.

\screenshotplaceholder{wazuh_local_rules_powershell.png}{Ausschnitt der eigenen PowerShell-Regeln in der Wazuh-Konfiguration}{fig:wazuh-local-rules-powershell}

\screenshotplaceholder{wazuh_rule_test_logtest.png}{Test einer eigenen Wazuh-Regel mit wazuh-logtest}{fig:wazuh-rule-test-logtest}

\begin{lstlisting}[language=XML,caption={Schematisches Beispiel einer Wazuh-Regel für PowerShell-Indikatoren}]
<group name="windows,powershell,scriptblock,">
  <rule id="100200" level="8">
    <if_group>windows</if_group>
    <field name="win.system.providerName">Microsoft-Windows-PowerShell</field>
    <match>EncodedCommand|FromBase64String|Invoke-Expression|DownloadString</match>
    <description>PowerShell ScriptBlock contains suspicious lexical indicator</description>
  </rule>
</group>
\end{lstlisting}

Das Listing dient als strukturelles Beispiel. Die finale Regel muss an die tatsächliche Feldstruktur der eingehenden Wazuh-Events angepasst und im Labor getestet werden.

\section{Verwendete Technologien und Artefakte}
Die folgenden Werkzeuge und Artefakte bilden die praktische Grundlage der Implementierung und der Auswertung.

\begin{table}[H]
\centering
\caption{Verwendete Technologien und Artefakte}
\small
\begin{tabularx}{\textwidth}{|L{5cm}|Y|}
\hline
\textbf{Technologie / Artefakt} & \textbf{Verwendung in der Arbeit}\\
\hline
\texttt{pipeline.py} & Hauptprogramm für Datensatzimport, Feature-Extraktion, lexikalische Klassifikation, Wazuh-Alert-Import und Metrikberechnung.\\
\hline
\codebreak{run\_dataset}{\_samples.ps1} & PowerShell-Runner zur Ausführung beziehungsweise ungefährlichen Protokollierung der Datensatz-Samples im Windows-Labor.\\
\hline
ScriptBlock-Logging & Windows-Telemetriequelle für die Erfassung der PowerShell-ScriptBlocks.\\
\hline
Wazuh Agent / Manager / Dashboard & Sammlung, Regelanalyse, Alarmierung und Export der PowerShell-Ereignisse.\\
\hline
Datensatz JSON v4 & Kuratierter Datensatz mit 300 gelabelten ScriptBlock-Samples.\\
\hline
Ergebnisdateien CSV/JSON & \texttt{sample\_results.csv}, \texttt{feature\_values.csv}, \texttt{metrics\_overall.csv}, \texttt{metrics\_by\_obfuscation.csv}, \texttt{metrics\_by\_dataset\_group.csv} und \texttt{dataset\_with\_results.json}.\\
\hline
\texttt{threshold\_tuning.csv} & Ergebnisdatei zur Auswahl des Schwellenwerts der lexikalischen Pipeline.\\
\hline
\end{tabularx}
\end{table}

\section{Python-Pipeline}
Die Python-Pipeline wurde als eigenständiges Auswertungspaket \texttt{thesis\_wazuh\_pipeline} implementiert. Sie verwendet ausschließlich Standardbibliotheken von Python und kann daher ohne zusätzliche Abhängigkeiten auf dem Analysehost ausgeführt werden. Die Pipeline übernimmt vier Aufgaben: Einlesen des JSON-Datensatzes, Extraktion lexikalischer Merkmale, scorebasierte Klassifikation und Berechnung der Vergleichsmetriken. Optional importiert sie exportierte Wazuh-Alerts und ordnet diese über die \texttt{WAZUH\_SAMPLE}-Marker den einzelnen Samples zu.

Der typische Aufruf für die lexikalische Auswertung lautet:
\begin{lstlisting}[language=bash,caption={Aufruf der lexikalischen Pipeline}]
python thesis_wazuh_pipeline/pipeline.py run \
  --dataset powershell_scriptblock_samples_final_thesis_v4.json \
  --out-dir thesis_wazuh_pipeline/results/lexical_only \
  --threshold 6
\end{lstlisting}

\screenshotplaceholder{pipeline_terminal_run.png}{Terminalausgabe eines Pipeline-Laufs mit Datensatzpfad, Schwellenwert und Ergebnisverzeichnis}{fig:pipeline-terminal-run}

Für die spätere Wazuh-Baseline wird zusätzlich ein Alert-Export übergeben:
\begin{lstlisting}[language=bash,caption={Aufruf der Pipeline mit Wazuh-Alert-Export}]
python thesis_wazuh_pipeline/pipeline.py run \
  --dataset powershell_scriptblock_samples_final_thesis_v4.json \
  --alerts wazuh_exports/alerts.json \
  --out-dir thesis_wazuh_pipeline/results/wazuh_standard_baseline \
  --threshold 6
\end{lstlisting}

\subsection{Feature-Extraktion}
Die Pipeline berechnet pro ScriptBlock einfache, erklärbare Merkmale. Dazu gehören Länge, Tokenanzahl, Entropie, Sonderzeichendichte, Ziffernanteil, Großbuchstabenanteil, Base64-Kandidaten, Backticks, Konkatenationsoperatoren, URLs, IP-ähnliche Muster, verdächtige PowerShell-Tokens, Alias-Tokens, LOLBin-Tokens und Muster für aufgeteilte Indikatoren.

\begin{lstlisting}[language=Python,caption={Vereinfachte Feature-Extraktion der Pipeline}]
def extract_features(scriptblock: str) -> dict:
    length = len(scriptblock)
    special_chars = len(re.findall(r"[^a-zA-Z0-9\s]", scriptblock))
    base64_candidates = BASE64_RE.findall(scriptblock)
    suspicious_count, suspicious_hits = count_token_hits(
        scriptblock, SUSPICIOUS_TOKENS
    )
    return {
        "length": length,
        "special_char_ratio": special_chars / max(length, 1),
        "entropy": shannon_entropy(scriptblock),
        "base64_candidate_count": len(base64_candidates),
        "backtick_count": scriptblock.count("`"),
        "concat_operator_count": scriptblock.count("+"),
        "suspicious_token_count": suspicious_count,
        "suspicious_tokens": suspicious_hits,
    }
\end{lstlisting}

\subsection{Schwellenwertlogik}
Die Klassifikation erfolgt über ein transparentes Punktemodell. Verdächtige Tokens, Hochrisiko-Tokens, Base64-Kandidaten, hohe Entropie, auffällige Sonderzeichendichte, Backticks, String-Konkatenation, Alias-Tokens und LOLBin-Indikatoren erhöhen den Score. Reine administrative Muster ohne verdächtige Indikatoren senken den Score leicht. In einer Schwellenwertanalyse auf dem finalen Datensatz erzielte der Schwellenwert 6 den besten F1-Score bei einer FPR unter 10 Prozent.

\begin{lstlisting}[language=Python,caption={Vereinfachte scorebasierte Klassifikation}]
def predict_lexical(scriptblock: str, threshold: int = 6) -> int:
    features = extract_features(scriptblock)
    score = 0
    if features["suspicious_token_count"] >= 1:
        score += 3
    if features["base64_candidate_count"] >= 1:
        score += 2
    if features["entropy"] >= 4.5:
        score += 1
    if features["backtick_count"] >= 2:
        score += 2
    if features["concat_operator_count"] >= 3:
        score += 2
    return int(score >= threshold)
\end{lstlisting}

\section{Metrikberechnung}
Die Metriken werden aus True Positives, False Positives, True Negatives und False Negatives berechnet. Für jede Vergleichsgruppe werden dieselben Funktionen genutzt, damit Wazuh, lexikalische Pipeline und kombinierter Ansatz vergleichbar bleiben.

\begin{lstlisting}[language=Python,caption={Berechnung zentraler Evaluierungsmetriken}]
def compute_metrics(y_true: list[int], y_pred: list[int]) -> dict:
    tp = sum(1 for t, p in zip(y_true, y_pred) if t == 1 and p == 1)
    fp = sum(1 for t, p in zip(y_true, y_pred) if t == 0 and p == 1)
    tn = sum(1 for t, p in zip(y_true, y_pred) if t == 0 and p == 0)
    fn = sum(1 for t, p in zip(y_true, y_pred) if t == 1 and p == 0)

    precision = tp / (tp + fp) if (tp + fp) else 0.0
    recall = tp / (tp + fn) if (tp + fn) else 0.0
    f1 = 2 * precision * recall / (precision + recall) if (precision + recall) else 0.0
    fpr = fp / (fp + tn) if (fp + tn) else 0.0

    return {"TP": tp, "FP": fp, "TN": tn, "FN": fn,
            "precision": precision, "recall": recall,
            "f1": f1, "false_positive_rate": fpr}
\end{lstlisting}

\section{Export der Ergebnisse}
Die Pipeline erzeugt maschinenlesbare Ergebnisdateien im CSV- und JSON-Format. Dazu gehören \texttt{sample\_results.csv}, \texttt{dataset\_with\_results.json}, \texttt{feature\_values.csv}, \texttt{wazuh\_alerts\_mapped.csv}, \texttt{metrics\_overall.csv}, \texttt{metrics\_by\_obfuscation.csv}, \texttt{metrics\_by\_dataset\_group.csv} und \texttt{metrics\_by\_tactic.csv}. Diese Dateien dienen als Grundlage für die Tabellen im Evaluationskapitel.

\screenshotplaceholder{python_pipeline_metrics_output.png}{Ausgabe der Python-Pipeline mit berechneten Metriken}{fig:python-pipeline-metrics-output}

\screenshotplaceholder{generated_metrics_files.png}{Erzeugte Ergebnisdateien der Pipeline im Ausgabeordner}{fig:generated-metrics-files}

\chapter{Evaluation}

\section{Zielsetzung der Evaluation}
Die Evaluation soll beantworten, ob standardmäßige Wazuh-Regeln obfuskierte PowerShell-Aktivität erkennen und ob lexikalische Merkmale die Erkennung ergänzen können. Bewertet werden nicht nur erkannte simuliert maliziöse Samples, sondern auch Fehlalarme bei normalen und schwerer abzugrenzenden benignen PowerShell-Samples.

\section{Versuchsaufbau}
Jeder Testfall wird in der isolierten Windows-11-Umgebung ausgeführt beziehungsweise ungefährlich protokolliert. Danach werden die erzeugten ScriptBlock-Logs und Wazuh-Alerts exportiert. Die Python-Pipeline verarbeitet dieselben ScriptBlocks und erzeugt eine zweite Entscheidung. Anschließend werden Wazuh, lexikalische Analyse und Kombination miteinander verglichen.

\screenshotplaceholder{experiment_dataset_overview.png}{Übersicht des kuratierten Datensatzes mit Labels und Obfuskierungsarten}{fig:experiment-dataset-overview}

\section{Versuchsablauf}
Der Ablauf pro Testfall ist standardisiert:
\begin{enumerate}
    \item Auswahl eines Datensatz-Eintrags.
    \item Ausführung oder ungefährliche Protokollierung des PowerShell-ScriptBlocks.
    \item Kontrolle des Windows-Eventlogs auf ScriptBlock-Ereignisse.
    \item Kontrolle und Export der Wazuh-Alerts.
    \item Zuordnung der Alerts über den Marker \texttt{WAZUH\_SAMPLE <ID>}.
    \item Berechnung der lexikalischen Features und Vorhersage.
    \item Berechnung der Metriken für Wazuh, lexikalische Pipeline und Kombination.
\end{enumerate}

\section{Metriken}
\begin{table}[H]
\centering
\caption{Metriken der Evaluation}
\begin{tabularx}{\textwidth}{|L{2.6cm}|Y|}
\hline
\textbf{Metrik} & \textbf{Bedeutung}\\
\hline
True Positive & Simuliert maliziöser ScriptBlock wird korrekt erkannt.\\
\hline
False Positive & Legitimer ScriptBlock wird fälschlicherweise als maliziös erkannt.\\
\hline
True Negative & Legitimer ScriptBlock wird korrekt nicht alarmiert.\\
\hline
False Negative & Simuliert maliziöser ScriptBlock wird nicht erkannt.\\
\hline
Precision & Anteil korrekter Alarme an allen Alarmen.\\
\hline
Recall & Anteil erkannter positiver Fälle an allen tatsächlich positiven Fällen.\\
\hline
F1-Score & Harmonisches Mittel aus Precision und Recall.\\
\hline
FPR & Anteil fälschlicher Alarme unter legitimen Fällen.\\
\hline
\end{tabularx}
\end{table}

\section{Schwellenwertanalyse}
Der Schwellenwert bestimmt, ab welchem Score ein ScriptBlock durch die lexikalische Pipeline als verdächtig klassifiziert wird. Ein niedriger Schwellenwert erhöht in der Regel den Recall, führt aber auch zu mehr False Positives. Ein höherer Schwellenwert reduziert Fehlalarme, übersieht jedoch mehr simuliert maliziöse Samples. Die folgende Tabelle basiert auf der Datei \texttt{threshold\_tuning.csv}. Für die weitere Auswertung wird der Schwellenwert 6 verwendet, da er den besten F1-Score bei einer FPR unter 10 Prozent erreicht.

\begin{table}[H]
\centering
\caption{Schwellenwertanalyse der lexikalischen Pipeline}
\begingroup
\small
\setlength{\tabcolsep}{3pt}
\begin{tabular}{|c|c|c|c|c|c|c|c|c|}
\hline
\textbf{Schw.} & \textbf{TP} & \textbf{FP} & \textbf{TN} & \textbf{FN} & \textbf{Prec.} & \textbf{Rec.} & \textbf{F1} & \textbf{FPR}\\
\hline
1 & 150 & 150 & 0 & 0 & 0{,}5000 & 1{,}0000 & 0{,}6667 & 1{,}0000\\
\hline
2 & 150 & 103 & 47 & 0 & 0{,}5929 & 1{,}0000 & 0{,}7444 & 0{,}6867\\
\hline
3 & 147 & 101 & 49 & 3 & 0{,}5927 & 0{,}9800 & 0{,}7387 & 0{,}6733\\
\hline
4 & 132 & 47 & 103 & 18 & 0{,}7374 & 0{,}8800 & 0{,}8024 & 0{,}3133\\
\hline
5 & 122 & 19 & 131 & 28 & 0{,}8652 & 0{,}8133 & 0{,}8385 & 0{,}1267\\
\hline
6 & 119 & 13 & 137 & 31 & 0{,}9015 & 0{,}7933 & 0{,}8440 & 0{,}0867\\
\hline
7 & 107 & 6 & 144 & 43 & 0{,}9469 & 0{,}7133 & 0{,}8137 & 0{,}0400\\
\hline
8 & 92 & 6 & 144 & 58 & 0{,}9388 & 0{,}6133 & 0{,}7419 & 0{,}0400\\
\hline
9 & 81 & 4 & 146 & 69 & 0{,}9529 & 0{,}5400 & 0{,}6894 & 0{,}0267\\
\hline
10 & 55 & 2 & 148 & 95 & 0{,}9649 & 0{,}3667 & 0{,}5314 & 0{,}0133\\
\hline
11 & 44 & 2 & 148 & 106 & 0{,}9565 & 0{,}2933 & 0{,}4490 & 0{,}0133\\
\hline
12 & 30 & 2 & 148 & 120 & 0{,}9375 & 0{,}2000 & 0{,}3297 & 0{,}0133\\
\hline
13 & 12 & 0 & 150 & 138 & 1{,}0000 & 0{,}0800 & 0{,}1481 & 0{,}0000\\
\hline
14 & 9 & 0 & 150 & 141 & 1{,}0000 & 0{,}0600 & 0{,}1132 & 0{,}0000\\
\hline
\end{tabular}
\endgroup
\end{table}

\screenshotplaceholder{threshold_tuning_chart.png}{Diagramm der Schwellenwertanalyse mit Precision, Recall, F1-Score und FPR}{fig:threshold-tuning-chart}

\section{Ergebnisdarstellung für Wazuh}
Für den finalen Wazuh-Lauf wurden alle 300 Samples des Datensatzes mit dem PowerShell-Runner in der Laborumgebung ausgeführt beziehungsweise protokolliert. Der Windows-Agent übertrug die PowerShell-Ereignisse an den Wazuh-Manager in WSL Kali. Aus der Datei \texttt{/var/ossec/logs/alerts/alerts.json} wurden 69 relevante Alert-Zeilen mit Datensatzbezug exportiert. Diese Alerts konnten 64 eindeutigen Samples zugeordnet werden.

Die Zuordnung erfolgte über die im ScriptBlock enthaltenen Sample-IDs wie \texttt{BEN-001}, \texttt{SIM-001} und \texttt{OBF-001}. Da Wazuh den ScriptBlock-Text vor der PowerShell-\allowbreak Stringformatierung protokolliert, erscheint die Markierung teilweise als Formatstring. Die Sample-ID selbst bleibt jedoch im ScriptBlock enthalten und kann durch die Pipeline zuverlässig extrahiert werden.

\begin{table}[H]
\centering
\caption{Status der Wazuh-Standardregel-Auswertung}
\begin{tabularx}{\textwidth}{|L{5.6cm}|Y|}
\hline
\textbf{Teilaufgabe} & \textbf{Status}\\
\hline
ScriptBlock-Samples mit eindeutiger Sample-ID & umgesetzt\\
\hline
PowerShell-Runner für Windows-Laborumgebung & umgesetzt\\
\hline
Import von Wazuh-Alerts im JSONL-Format & umgesetzt\\
\hline
Zuordnung von Alerts zu Sample-IDs & umgesetzt: 69 Alert-Zeilen, 64 eindeutige Samples\\
\hline
Gemessene Wazuh-Standardregelwerte & umgesetzt\\
\hline
\end{tabularx}
\end{table}

\begin{table}[H]
\centering
\caption{Gemessene Wazuh-Baseline-Ergebnisse}
\begingroup
\small
\setlength{\tabcolsep}{3pt}
\begin{tabular}{|L{4cm}|c|c|c|c|c|c|c|c|}
\hline
\textbf{Methode} & \textbf{TP} & \textbf{FP} & \textbf{TN} & \textbf{FN} & \textbf{Prec.} & \textbf{Rec.} & \textbf{F1} & \textbf{FPR}\\
\hline
Wazuh-Standardregeln & 30 & 34 & 116 & 120 & 0{,}4688 & 0{,}2000 & 0{,}2804 & 0{,}2267\\
\hline
\end{tabular}
\endgroup
\end{table}

\begin{table}[H]
\centering
\caption{Ausgelöste Wazuh-Regeln im Alert-Export}
\small
\begin{tabularx}{\textwidth}{|L{2cm}|C{1.8cm}|Y|}
\hline
\textbf{Regel-ID} & \textbf{Alerts} & \textbf{Beschreibung}\\
\hline
\texttt{91809} & 29 & Powershell script may be using Base64 decoding method\\
\hline
\texttt{91816} & 16 & Powershell script querying system environment variables\\
\hline
\texttt{91815} & 13 & Powershell executing process discovery\\
\hline
\texttt{91837} & 5 & Powershell executed \enquote{Get-Content -Stream or Invoke-Expresion}. Possible string execution as code\\
\hline
\texttt{91819} & 4 & Powershell script searching filesystem\\
\hline
\texttt{91823} & 2 & Powershell script used \enquote{Invoke-command} cmdlet to execute code on remote computer\\
\hline
\end{tabularx}
\end{table}

Die Ergebnisse zeigen, dass die Wazuh-Standardregeln in dieser Konfiguration vor allem Base64-Decoding zuverlässig erkennen. Andere Obfuskierungsarten werden deutlich seltener erkannt. Gleichzeitig erzeugen generische PowerShell-Regeln Fehlalarme bei legitimen administrativen Samples, etwa bei Prozessabfragen, Dateisystemsuche oder Umgebungsvariablenabfragen.

\begin{table}[H]
\centering
\caption{Wazuh-Erkennung nach Obfuskierungsart}
\small
\begin{tabularx}{\textwidth}{|L{3.6cm}|C{1.1cm}|C{1.1cm}|C{1.5cm}|Y|}
\hline
\textbf{Obfuskierungsart} & \textbf{TP} & \textbf{FN} & \textbf{Recall} & \textbf{Interpretation}\\
\hline
Base64-Indikatoren & 20 & 0 & 1{,}0000 & Vollständig durch Regel \texttt{91809} erkannt.\\
\hline
String-Konkatenation & 2 & 18 & 0{,}1000 & Meist nicht durch Standardregeln abgedeckt.\\
\hline
Backticks & 2 & 18 & 0{,}1000 & Einfache Token-Unterbrechung reduziert die Wazuh-Erkennung deutlich.\\
\hline
Gemischte Schreibweise & 2 & 18 & 0{,}1000 & Groß- und Kleinschreibung wird nur begrenzt robust erkannt.\\
\hline
Alias-/Befehlsrekonstruktion & 2 & 18 & 0{,}1000 & Rekonstruierte Befehle fallen selten unter vorhandene Regeln.\\
\hline
Kombinierte Obfuskation & 1 & 19 & 0{,}0500 & Mehrschichtige Obfuskation wird durch Standardregeln kaum erkannt.\\
\hline
\end{tabularx}
\end{table}

\screenshotplaceholder{wazuh_alert_view_export.png}{Wazuh-Dashboard mit PowerShell-Alerts und Exportfunktion}{fig:wazuh-alert-view-export}
\screenshotplaceholder{wazuh_alert_export_file.png}{Exportierte Wazuh-Alerts im JSONL-Format}{fig:wazuh-alert-export-file}

\section{Ergebnisdarstellung für lexikalische Analyse}
Die lexikalische Pipeline wurde auf dem finalen Datensatz mit einem Schwellenwert von 6 ausgeführt. Dieser Schwellenwert wurde gewählt, weil er in der Schwellenwertanalyse den besten F1-Score bei einer FPR unter 10 Prozent erreichte.

\begin{table}[H]
\centering
\caption{Gesamtergebnis der lexikalischen Pipeline}
\begingroup
\small
\setlength{\tabcolsep}{3pt}
\begin{tabular}{|L{3.5cm}|c|c|c|c|c|c|c|c|}
\hline
\textbf{Methode} & \textbf{TP} & \textbf{FP} & \textbf{TN} & \textbf{FN} & \textbf{Prec.} & \textbf{Rec.} & \textbf{F1} & \textbf{FPR}\\
\hline
Lexikalische Pipeline & 119 & 13 & 137 & 31 & 0{,}9015 & 0{,}7933 & 0{,}8440 & 0{,}0867\\
\hline
\end{tabular}
\endgroup
\end{table}

\begin{table}[H]
\centering
\caption{Lexikalische Ergebnisse nach Datensatzgruppe}
\begingroup
\small
\setlength{\tabcolsep}{3pt}
\begin{tabular}{|L{4.4cm}|c|c|c|c|c|c|c|c|}
\hline
\textbf{Gruppe} & \textbf{TP} & \textbf{FP} & \textbf{TN} & \textbf{FN} & \textbf{Prec.} & \textbf{Rec.} & \textbf{F1} & \textbf{FPR}\\
\hline
Direkte Indikatoren & 18 & 0 & 0 & 12 & 1{,}0000 & 0{,}6000 & 0{,}7500 & 0{,}0000\\
\hline
Obfuskierte Indikatoren & 101 & 0 & 0 & 19 & 1{,}0000 & 0{,}8417 & 0{,}9140 & 0{,}0000\\
\hline
Normale benigne Samples & 0 & 2 & 98 & 0 & 0{,}0000 & 0{,}0000 & 0{,}0000 & 0{,}0200\\
\hline
Schwere benigne Samples & 0 & 11 & 39 & 0 & 0{,}0000 & 0{,}0000 & 0{,}0000 & 0{,}2200\\
\hline
\end{tabular}
\endgroup
\end{table}

\section{Vergleich der drei Erkennungsgruppen}
Nach dem Import der Wazuh-Alerts können alle drei Vergleichsgruppen direkt gegenübergestellt werden. Wazuh allein erreicht in dieser Konfiguration einen niedrigen Recall, weil viele obfuskierte Varianten keine Standardregel auslösen. Die lexikalische Pipeline erzielt den höchsten F1-Score und zugleich eine deutlich niedrigere FPR. Die Kombination nach dem ODER-Prinzip erkennt ein zusätzliches simuliert maliziöses Sample gegenüber der lexikalischen Pipeline, erhöht aber die Zahl der False Positives von 13 auf 39.

\screenshotplaceholder{metrics_comparison_chart.png}{Diagramm zum Vergleich von Precision, Recall, F1-Score und FPR}{fig:metrics-comparison-chart}

\begin{table}[H]
\centering
\caption{Gesamtvergleich der Erkennungsgruppen}
\begingroup
\small
\setlength{\tabcolsep}{3pt}
\begin{tabular}{|L{4cm}|c|c|c|c|c|c|c|c|}
\hline
\textbf{Methode} & \textbf{TP} & \textbf{FP} & \textbf{TN} & \textbf{FN} & \textbf{Prec.} & \textbf{Rec.} & \textbf{F1} & \textbf{FPR}\\
\hline
Wazuh-Standardregeln & 30 & 34 & 116 & 120 & 0{,}4688 & 0{,}2000 & 0{,}2804 & 0{,}2267\\
\hline
Lexikalische Pipeline & 119 & 13 & 137 & 31 & 0{,}9015 & 0{,}7933 & 0{,}8440 & 0{,}0867\\
\hline
Kombination & 120 & 39 & 111 & 30 & 0{,}7547 & 0{,}8000 & 0{,}7767 & 0{,}2600\\
\hline
\end{tabular}
\endgroup
\end{table}
\section{Auswertung nach Obfuskierungsart}
Die Ergebnisse zeigen deutliche Unterschiede zwischen den Obfuskierungsarten. Base64-Indikatoren wurden durch die lexikalische Pipeline vollständig erkannt. String-Konkatenation wurde ebenfalls sehr gut erkannt. Backticks und gemischte Groß- und Kleinschreibung wurden überwiegend erkannt, während Alias- beziehungsweise Befehlsrekonstruktion und kombinierte Obfuskation schwieriger waren.

\begin{table}[H]
\centering
\caption{Lexikalische Ergebnisse nach Obfuskierungsart}
\begingroup
\small
\setlength{\tabcolsep}{3pt}
\begin{tabular}{|L{3.8cm}|c|c|c|c|c|c|c|c|}
\hline
\textbf{Obfuskierungsart} & \textbf{TP} & \textbf{FP} & \textbf{TN} & \textbf{FN} & \textbf{Prec.} & \textbf{Rec.} & \textbf{F1} & \textbf{FPR}\\
\hline
Base64 & 20 & 0 & 0 & 0 & 1{,}0000 & 1{,}0000 & 1{,}0000 & 0{,}0000\\
\hline
String-Konkatenation & 19 & 0 & 0 & 1 & 1{,}0000 & 0{,}9500 & 0{,}9744 & 0{,}0000\\
\hline
Backticks & 17 & 0 & 0 & 3 & 1{,}0000 & 0{,}8500 & 0{,}9189 & 0{,}0000\\
\hline
Gemischte Schreibweise & 17 & 0 & 0 & 3 & 1{,}0000 & 0{,}8500 & 0{,}9189 & 0{,}0000\\
\hline
Alias/Rekonstruktion & 14 & 0 & 0 & 6 & 1{,}0000 & 0{,}7000 & 0{,}8235 & 0{,}0000\\
\hline
Kombinierte Obfuskation & 14 & 0 & 0 & 6 & 1{,}0000 & 0{,}7000 & 0{,}8235 & 0{,}0000\\
\hline
Keine Obfuskation & 18 & 13 & 137 & 12 & 0{,}5806 & 0{,}6000 & 0{,}5902 & 0{,}0867\\
\hline
\end{tabular}
\endgroup
\end{table}

\section{Bewertung der False Positives}
Die FPR der lexikalischen Pipeline beträgt insgesamt 0{,}0867. Von 150 benignen Samples wurden 13 fälschlicherweise als verdächtig klassifiziert. Dabei stammen 11 Fehlalarme aus der Gruppe der schweren benignen Samples und 2 aus der Gruppe der normalen benignen Samples. Dies zeigt, dass die bewusst schwierigeren benignen Fälle für die Bewertung der praktischen Nutzbarkeit wichtig sind. Ohne diese Gruppe würde die FPR der lexikalischen Pipeline zu optimistisch erscheinen.

\chapter{Diskussion}

\section{Interpretation der Ergebnisse}
Die Diskussion bewertet, ob die lexikalische Analyse einen echten Mehrwert gegenüber Wazuh liefert. Ein Mehrwert liegt vor, wenn zusätzliche simulierte maliziöse oder obfuskierte ScriptBlocks erkannt werden, ohne dass die FPR unverhältnismäßig steigt. Dabei muss zwischen Recall-Gewinn und Alarmqualität unterschieden werden.

\section{Grenzen des Datensatzes}
Der Datensatz ist kontrolliert und dadurch reproduzierbar, aber nicht vollständig repräsentativ für alle realen PowerShell-Angriffe. Die Auswahl der legitimen Skripte, der simulierten Angriffsmuster und der Obfuskierungstechniken beeinflusst die Ergebnisse. Deshalb wird transparent dokumentiert, wie die Samples entstanden sind und welche Einschränkungen gelten.

\section{Grenzen von ScriptBlock-Logging}
ScriptBlock-Logging ist eine wertvolle Datenquelle, aber nicht unfehlbar. Die Erkennung hängt davon ab, dass Logging korrekt aktiviert ist, Ereignisse nicht verloren gehen und Wazuh die relevanten Felder vollständig erhält. Außerdem liefert ScriptBlock-Logging zwar Inhaltsinformationen, aber keine vollständige Kontextanalyse aller Systemaktivitäten.

\section{Grenzen regelbasierter Erkennung}
Regelbasierte Erkennung ist erklärbar, kann aber durch veränderte Muster umgangen werden. Je spezifischer eine Regel ist, desto geringer ist meist die FPR, aber desto höher kann das Risiko von False Negatives sein. Je allgemeiner eine Regel ist, desto mehr Varianten erkennt sie, aber desto eher entstehen Fehlalarme.

\section{Nutzen lexikalischer Merkmale}
Lexikalische Merkmale können helfen, verdächtige Strukturmerkmale zu erkennen, die nicht an ein einzelnes Klartexttoken gebunden sind. Dazu gehören hohe Entropie, ungewöhnliche Sonderzeichendichte oder viele Backticks. Gleichzeitig sind solche Merkmale nur Indikatoren. Sie beweisen keine Malizität und müssen im Kontext interpretiert werden.

\section{Bedrohungen der Validität}
Die Validität der Arbeit wird durch mehrere Faktoren beeinflusst:
\begin{itemize}
    \item begrenzte Größe des Datensatzes,
    \item künstlich erzeugte simulierte Angriffsmuster,
    \item begrenzte Auswahl an Obfuskierungstechniken,
    \item mögliche Unterschiede zwischen Laborumgebung und produktiven Umgebungen,
    \item manuelle Entscheidungen bei Labeling und Schwellenwertwahl,
    \item mögliche Verzerrung durch die Auswahl legitimer Administrationsskripte.
\end{itemize}

Diese Punkte werden offen dokumentiert, damit die Ergebnisse realistisch eingeordnet werden können.

\section{Praktische Empfehlungen}
Aus den gemessenen Ergebnissen können praktische Empfehlungen abgeleitet werden. Dazu zählen die Aktivierung von ScriptBlock-Logging, die zentrale Sammlung von PowerShell-Ereignissen, die Ergänzung eigener Wazuh-Regeln für organisationsspezifische Muster und die Nutzung einfacher Zusatzmerkmale zur Priorisierung auffälliger ScriptBlocks. Für produktive Umgebungen wird eine solche Erkennung jedoch durch weitere Telemetriequellen ergänzt, etwa Prozessstartdaten, Sysmon-Ereignisse oder Netzwerkmetadaten.

\chapter{Herausforderungen während der Umsetzung}

\section{Wazuh-Feldzuordnung}
Eine zentrale Herausforderung besteht darin, dass Wazuh-Events je nach Agent-, Manager- und Windows-Konfiguration unterschiedliche Feldnamen und verschachtelte Strukturen enthalten können. Für die Auswertung ist deshalb entscheidend, die tatsächlichen Felder der eingehenden PowerShell-Ereignisse zu prüfen und die Regelbedingungen sowie den Alert-Import der Pipeline daran anzupassen.

\section{Weiterleitung der ScriptBlock-Logs}
ScriptBlock-Logging muss nicht nur lokal aktiviert sein, sondern die relevanten Ereignisse müssen auch vom Wazuh-Agenten gesammelt und an den Manager übertragen werden. Fehlende Ereignisse können sonst wie eine schlechte Detektionsleistung wirken, obwohl tatsächlich die TelemetrieÜbertragung unvollständig ist.

\section{Zuordnung von Alerts zu Samples}
Für den Vergleich müssen Wazuh-Alerts eindeutig auf Datensatz-Samples zurückgeführt werden. Deshalb enthält jeder ScriptBlock eine \texttt{WAZUH\_SAMPLE}-Markierung. Trotzdem muss geprüft werden, ob Wazuh den vollständigen ScriptBlock-Inhalt übernimmt oder ob Felder gekürzt, normalisiert oder auf mehrere Events verteilt werden.

\section{False Positives bei schwer benignen Samples}
Die schwer benignen Samples erzeugen bewusst höhere lexikalische Auffälligkeit, etwa durch längere Strings, sichere Base64-Inhalte oder komplexere administrative Verarbeitung. Diese Fälle sind wichtig, weil sie zeigen, ob ein Ansatz nur einfache Beispiele erkennt oder auch bei realistischeren administrativen Skripten eine akzeptable FPR behält.

\section{Sicheres Design maliziös wirkender Muster}
Die simuliert maliziösen Samples enthalten verdächtige Begriffe und obfuskierte Strukturen, führen jedoch keine Schadfunktion aus. Dadurch bleibt der Laborversuch ethisch und technisch kontrollierbar. Gleichzeitig begrenzt diese Sicherheitsentscheidung die Übertragbarkeit auf reale Angriffe, was in der Diskussion berücksichtigt werden muss.

\chapter{Fazit und Ausblick}

\section{Zusammenfassung}
Diese Arbeit entwickelt einen reproduzierbaren Evaluierungsaufbau zur Untersuchung obfuskierter PowerShell-Aktivität unter Windows 11. Der finale Datensatz umfasst 300 ScriptBlock-Samples mit ausgeglichener Verteilung zwischen benignen und simuliert maliziösen Fällen. Die simuliert maliziösen Samples enthalten direkte und obfuskierte Indikatoren, bleiben aber ungefährlich und nicht-ausführend. Zusätzlich wurde eine Python-Pipeline implementiert, die lexikalische Merkmale extrahiert, einen transparenten Score berechnet, Wazuh-Alert-Exporte einliest und zentrale Metriken erzeugt.

Im finalen Laborlauf wurden alle Samples ausgeführt beziehungsweise protokolliert. Der Windows-Wazuh-Agent übertrug PowerShell-Ereignisse an den Wazuh-Manager in WSL Kali. Aus dem Wazuh-Alertlog wurden 69 relevante Alert-Zeilen exportiert und 64 eindeutigen Samples zugeordnet. Damit konnte die zuvor offene Wazuh-Baseline quantitativ bewertet werden.

\section{Bewertung der Zielerreichung}
Die zentralen Umsetzungsziele wurden erreicht: Der kuratierte Datensatz wurde erstellt, die ungefährlichen Testfälle enthalten eindeutige Sample-IDs, die Python-Pipeline extrahiert lexikalische Merkmale und berechnet die zentralen Metriken, die Schwellenwertanalyse begründet den verwendeten Schwellenwert 6, und die Wazuh-Standardregel-Baseline wurde anhand real exportierter Alerts ausgewertet.

\begin{table}[H]
\centering
\caption{Bewertung der Zielerreichung}
\small
\begin{tabularx}{\textwidth}{|L{4cm}|Y|L{2cm}|}
\hline
\textbf{Ziel} & \textbf{Bewertung} & \textbf{Status}\\
\hline
Kuratierter Datensatz & 300 gelabelte ScriptBlock-Samples mit benignen, schwer benignen, direkten und obfuskierten Indikatorfällen. & erreicht\\
\hline
Python-Pipeline & Feature-Extraktion, Schwellenwertentscheidung, Metrikberechnung und Alert-Import sind implementiert. & erreicht\\
\hline
Lexikalische Evaluation & Precision, Recall, F1-Score und FPR wurden für den finalen Datensatz berechnet. & erreicht\\
\hline
Wazuh-Baseline & 69 relevante Wazuh-Alert-Zeilen wurden importiert und 64 eindeutigen Samples zugeordnet. & erreicht\\
\hline
Hybridvergleich & Wazuh, lexikalische Pipeline und ODER-Kombination wurden vollständig verglichen. & erreicht\\
\hline
\end{tabularx}
\end{table}

\section{Beantwortung der Forschungsfragen}
Die Wazuh-Standardregeln erkannten 30 von 150 simuliert maliziösen Samples und erzeugten 34 Fehlalarme bei 150 benignen Samples. Daraus ergeben sich eine Precision von 0{,}4688, ein Recall von 0{,}2000, ein F1-Score von 0{,}2804 und eine FPR von 0{,}2267. Besonders stark war Wazuh bei Base64-Decoding, wo alle 20 entsprechenden obfuskierten Samples erkannt wurden. Deutlich schwächer war die Standardregel-Erkennung bei String-Konkatenation, Backticks, gemischter Schreibweise, Alias-/Befehlsrekonstruktion und kombinierter Obfuskation.

Die lexikalische Pipeline erkannte 119 von 150 simuliert maliziösen Samples und erzeugte 13 Fehlalarme. Damit erreicht sie eine Precision von 0{,}9015, einen Recall von 0{,}7933, einen F1-Score von 0{,}8440 und eine FPR von 0{,}0867. Die ODER-Kombination aus Wazuh und lexikalischer Pipeline erkannte 120 simuliert maliziöse Samples, erhöhte aber die False Positives auf 39. Dadurch stieg der Recall leicht auf 0{,}8000, während Precision und F1-Score gegenüber der lexikalischen Pipeline sanken. Für diesen Datensatz ist die lexikalische Pipeline daher der stärkste Einzelansatz; die einfache ODER-Kombination liefert nur einen kleinen Recall-Gewinn bei deutlich höherem Alarmaufkommen.

\section{Ausblick}
Auf Basis der Messergebnisse bieten sich mehrere Erweiterungen an. Erstens könnten eigene Wazuh-Regeln gezielt für die schwach erkannten Obfuskierungsarten ergänzt und getrennt von der Standardregel-Baseline bewertet werden. Zweitens könnte die Kombination nicht als einfache ODER-Verknüpfung umgesetzt werden, sondern als priorisierendes Modell, das Wazuh-Regeln und lexikalische Scores gewichtet. Drittens wären zusätzliche Telemetriequellen sinnvoll, etwa Sysmon-Prozessdaten, Prozessketten oder Netzwerkmetadaten. Weitere Arbeiten könnten außerdem AST-basierte Merkmale, zusätzliche PowerShell-Versionen, komplexere Obfuskierungstechniken oder maschinelle Lernverfahren auf Basis der bereits extrahierten Features untersuchen.

%--------------------------------------------------------------------
% Literaturverzeichnis
%--------------------------------------------------------------------
\begin{thebibliography}{99}

\bibitem{microsoft-logging-windows}
Microsoft. \emph{about\_Logging\_Windows}. Microsoft Learn, aktualisiert am 18.01.2026. Online verfügbar: \url{https://learn.microsoft.com/en-us/powershell/module/microsoft.powershell.core/about/about_logging_windows?view=powershell-7.6}; abgerufen am 24.05.2026.

\bibitem{microsoft-logging}
Microsoft. \emph{about\_Logging}. Microsoft Learn, aktualisiert am 29.09.2025. Online verfügbar: \url{https://learn.microsoft.com/en-us/powershell/module/microsoft.powershell.core/about/about_logging?view=powershell-5.1}; abgerufen am 24.05.2026.

\bibitem{wazuh-log-collection}
Wazuh Inc. \emph{Configuring log collection for different operating systems}. Wazuh Documentation. Online verfügbar: \url{https://documentation.wazuh.com/current/user-manual/capabilities/log-data-collection/configuration.html}; abgerufen am 24.05.2026.

\bibitem{wazuh-rules}
Wazuh Inc. \emph{Rules Syntax -- Ruleset XML syntax}. Wazuh Documentation. Online verfügbar: \url{https://documentation.wazuh.com/current/user-manual/ruleset/ruleset-xml-syntax/rules.html}; abgerufen am 24.05.2026.

\bibitem{wazuh-custom-rules}
Wazuh Inc. \emph{Custom rules}. Wazuh Documentation. Online verfügbar: \url{https://documentation.wazuh.com/current/user-manual/ruleset/rules/custom.html}; abgerufen am 24.05.2026.

\bibitem{mitre-powershell}
MITRE. \emph{Command and Scripting Interpreter: PowerShell, Technique T1059.001}. MITRE ATT\&CK. Online verfügbar: \url{https://attack.mitre.org/techniques/T1059/001/}; abgerufen am 24.05.2026.

\bibitem{nist-idps}
Scarfone, Karen; Mell, Peter. \emph{Guide to Intrusion Detection and Prevention Systems (IDPS)}. NIST Special Publication 800-94, 2007. Online verfügbar: \url{https://csrc.nist.gov/pubs/sp/800/94/final}; abgerufen am 24.05.2026.

\bibitem{shannon1948}
Shannon, Claude E. \emph{A Mathematical Theory of Communication}. The Bell System Technical Journal, Vol. 27, 1948, S. 379--423 und 623--656. DOI: \url{https://doi.org/10.1002/j.1538-7305.1948.tb01338.x}; abgerufen am 24.05.2026.

\bibitem{elastic-powershell-winlogbeat}
Elastic. \emph{PowerShell module -- Winlogbeat Reference}. Online verfügbar: \url{https://www.elastic.co/docs/reference/beats/winlogbeat/winlogbeat-module-powershell}; abgerufen am 24.05.2026.

\bibitem{splunk-scriptblock}
Splunk. \emph{Hunting for Malicious PowerShell using Script Block Logging}. Splunk Blog, 17.09.2021. Online verfügbar: \url{https://www.splunk.com/en_us/blog/security/hunting-for-malicious-powershell-using-script-block-logging.html}; abgerufen am 24.05.2026.

\bibitem{atomic-red-team}
Red Canary. \emph{Atomic Red Team: T1059.001 PowerShell}. Online verfügbar: \url{https://github.com/redcanaryco/atomic-red-team/blob/master/atomics/T1059.001/T1059.001.md}; abgerufen am 28.06.2026.

\bibitem{sigmahq-powershell}
SigmaHQ. \emph{Sigma rules for Windows PowerShell}. Online verfügbar: \url{https://github.com/SigmaHQ/sigma/tree/master/rules/windows/powershell}; abgerufen am 28.06.2026.

\bibitem{lolbas}
LOLBAS Project. \emph{Living Off The Land Binaries, Scripts and Libraries}. Online verfügbar: \url{https://lolbas-project.github.io/}; abgerufen am 28.06.2026.

\bibitem{invoke-obfuscation}
Bohannon, Daniel. \emph{Invoke-Obfuscation}. Online verfügbar: \url{https://github.com/danielbohannon/Invoke-Obfuscation}; abgerufen am 28.06.2026.

\bibitem{elastic-detection-rules}
Elastic. \emph{Detection rules for Windows}. Online verfügbar: \url{https://github.com/elastic/detection-rules/tree/main/rules/windows}; abgerufen am 28.06.2026.

\end{thebibliography}

%--------------------------------------------------------------------
% Anhang
%--------------------------------------------------------------------
\appendix

\chapter{Datensatzstruktur}
\begin{lstlisting}[language={},caption={JSON-Schema des finalen Datensatzes in vereinfachter Form}]
{
  "id": "BEN-001",
  "label": "benign",
  "class_id": 0,
  "dataset_group": "normal_benign",
  "scenario": "process_inventory_normal_v1",
  "mitre_attack_id": "N/A",
  "obfuscation_type": "none",
  "scriptblock": "$sampleId = 'BEN-001'; ... WAZUH_SAMPLE BEN-001 ...",
  "safety_class": "benign_safe_execution",
  "features": {},
  "wazuh_alert_observed": null,
  "wazuh_rule_id_observed": null,
  "lexical_score": null,
  "lexical_prediction": null
}
\end{lstlisting}

\chapter{Pipeline-Artefakte}
\begin{longtable}{|L{6.3cm}|L{9cm}|}
\hline
\textbf{Datei} & \textbf{Zweck}\\
\hline
\texttt{pipeline.py} & Hauptpipeline für Datensatzimport, Feature-Extraktion, lexikalische Klassifikation, Wazuh-Alert-Import und Metrikberechnung.\\
\hline
\codebreak{run\_dataset}{\_samples.ps1} & PowerShell-Runner zur Ausführung der Datensatz-Samples in der isolierten Windows-11-Laborumgebung.\\
\hline
\texttt{threshold\_tuning.csv} & Tabelle zur Auswahl des lexikalischen Schwellenwerts.\\
\hline
\texttt{metrics\_overall.csv} & Gesamtergebnisse pro Erkennungsmethode.\\
\hline
\codebreak{metrics\_by}{\_obfuscation.csv} & Ergebnisse pro Obfuskierungsart.\\
\hline
\codebreakthree{metrics\_by}{\_dataset}{\_group.csv} & Ergebnisse pro Datensatzgruppe.\\
\hline
\codebreakthree{dataset}{\_with}{\_results.json} & Datensatz mit ergänzten Feature-, Wazuh- und Pipeline-Ergebnisfeldern.\\
\hline
\end{longtable}

\chapter{Experimentstatus}
\begin{longtable}{|L{4cm}|L{10cm}|}
\hline
\textbf{Feld} & \textbf{Eintrag}\\
\hline
Datensatz & \codebreakthree{powershell\_scriptblock}{\_samples\_final}{\_thesis\_v4.json}\\
\hline
Anzahl Samples & 300\\
\hline
Benigne Samples & 150\\
\hline
Simuliert maliziöse Samples & 150\\
\hline
Obfuskierte Samples & 120\\
\hline
Lexikalischer Schwellenwert & 6\\
\hline
Lexikalische Precision & 0{,}9015\\
\hline
Lexikalischer Recall & 0{,}7933\\
\hline
Lexikalischer F1-Score & 0{,}8440\\
\hline
Lexikalische FPR & 0{,}0867\\
\hline
Wazuh-Alert-Export & 69 relevante Alert-Zeilen aus \texttt{/var/ossec/logs/alerts/alerts.json}; 64 eindeutige Sample-IDs zugeordnet.\\
\hline
Wazuh Precision & 0{,}4688\\
\hline
Wazuh Recall & 0{,}2000\\
\hline
Wazuh F1-Score & 0{,}2804\\
\hline
Wazuh FPR & 0{,}2267\\
\hline
Hybrid F1-Score & 0{,}7767\\
\hline
\end{longtable}

\end{document}
